% Source: source/普通高考/2016/2016全国2文(甘肃,青海,内蒙古,黑龙江,吉林,辽宁,海南,宁夏,新疆,西藏,陕西,重庆).pdf, page 1.
% Type: algorithm flowchart; the original PDF contains vector line art (no embedded bitmap).
% Grid evidence: tmp/review_2016_national2_liberal/grid/page1_grid100.jpg;
% comparison crop: tmp/review_2016_national2_liberal/crops/q9_original_tight.png.
% Node-edge list: 开始→输入x,n→k=0,s=0→输入a→(s=s·x+a,k=k+1)→k>n；
% k>n(是)→输出s→结束；k>n(否)→右侧回边→输入a入口的主干汇合点。
% Solid segments: all box borders and directed connectors. Dashed segments: none.
% Labels 是/否 are placed beside their corresponding decision exits. Every node
% uses local parindent=0pt and centered text; loop arrow ends at the merge point
% and does not create a second arrow on the main input-a edge.
\begin{tikzpicture}[
  x=1cm,
  y=0.9cm,
  >=Stealth,
  line width=0.45pt,
  flowtext/.style={font=\normalsize\setlength{\parindent}{0pt}\noindent,
    anchor=center,align=center},
  every node/.style={flowtext,inner sep=1pt},
  flow/.style={flowtext,draw,line width=0.45pt},
  startstop/.style={flow, rounded corners=3pt, minimum width=1.2cm, minimum height=0.70cm},
  process/.style={flow, minimum height=0.72cm, align=center},
  io/.style={flow, trapezium, trapezium left angle=74, trapezium right angle=106,
    minimum height=0.72cm, align=center},
  decision/.style={flow, diamond, aspect=2.8, minimum width=3.2cm,
    minimum height=1.0cm, inner sep=1pt, align=center}
]
  \node[startstop] (start) at (0,9.65) {开始};
  \node[io, minimum width=2.30cm] (input) at (0,8.25) {输入 $x,n$};
  \node[process, minimum width=2.45cm] (init) at (0,6.95) {$k=0, s=0$};
  \node[io, minimum width=2.10cm] (inputa) at (0,5.65) {输入 $a$};
  \node[process, minimum width=2.75cm, minimum height=1.05cm] (update) at (0,4.15)
    {$\begin{gathered}s=s\cdot x+a\cr k=k+1\end{gathered}$};
  \node[decision] (check) at (0,2.55) {$k>n$};
  \node[io, minimum width=1.75cm] (output) at (0,1.05) {输出 $s$};
  \node[startstop] (finish) at (0,-0.25) {结束};

  % Main downward path.
  \draw[->] (start.south) -- (input.north);
  \draw[->] (input.south) -- (init.north);
  \draw[->] (init.south) -- (inputa.north);
  \draw[->] (inputa.south) -- (update.north);
  \draw[->] (update.south) -- (check.north);
  \draw[->] (output.south) -- (finish.north);

  % 否: right-side return to the input-a entry; the left-pointing arrow ends
  % at the main vertical edge, while the main init→input-a arrow remains unique.
  \coordinate (loopRight) at (2.65,2.55);
  \coordinate (loopTop) at (2.65,6.22);
  \coordinate (loopJoin) at (0,6.22);
  \draw[->] (check.east) -- node[pos=0.28,above] {否} (loopRight) -- (loopTop) -- (loopJoin);
  \draw[->] (check.south) -- node[midway,right] {是} (output.north);
\end{tikzpicture}
